% GENERATED FILE. Edit canonical lessons, then run npm run generate:books.
% canonical-source-hash: fnv1a64:167b37d149edde4f
% canonical-lessons: FA-W15-alef, FA-W15-lam, FA-W15-sin, FA-W15-mim, FA-W15-joining, FA-W15-be, FA-W15-te-nun, FA-W15-he, FA-W15-vav, FA-C15-practice

\chapter{Nine Letters, and the Words You Already Knew}
\label{ch:fa-script-nine}
\hlchaptermodality{15}

\begin{chapteropening}
\textbf{By the end of this chapter:} \emph{I can write nine Perso-Arabic letters from memory, read salâm, bale, na, mamnun, nân, nâm and mâh off the page rather than from the romanization, and explain why Persian leaves its short vowels unwritten.}

Seven words the book had only ever given by ear, read letter by letter from the right -- and the reason a Persian word can be read only once you have heard it.
\end{chapteropening}

\section[alef]{\fa{ا} — one stroke, and the direction of the page}

\label{lesson:FA-W15-alef}



\begin{quote}
Return to the right-to-left direction and the one alef stroke you first copied beside \emph{salâm}. Now make both durable before adding a second letter.
\end{quote}

\subsection*{Script — first, which way the page runs}
English marches \textbf{left to right}. Persian marches \textbf{right to left}.

Your pen starts at the \textbf{right edge} of a word and travels \textbf{towards the left}. So in a Persian word the letter you meet \textbf{first} is the \textbf{rightmost} one, and the letter you meet \textbf{last} is the leftmost — the exact reverse of what your eye has been trained on for your whole reading life.

This takes a few minutes to stop feeling strange and then stops mattering entirely. Numbers, incidentally, still run left to right, which surprises people more than the letters do.

\subsection*{Script — the letter itself}
\begin{quote}
\fa{ا}
\end{quote}

\textbf{One downstroke.} That is the whole letter, and there is nothing in the alphabet simpler. Draw it from the top down, straight, and lift.

Its name is \textbf{alef}. What it \emph{sounds} like is a longer story than what it looks like, and this book will take that slowly: for now, treat alef as the letter that carries a long \textbf{â}, the vowel in English \emph{father}.

\subsection*{Your turn}
\begin{itemize}
\raggedright
  \item \emph{Write it:} \fa{ا} — one stroke, top to bottom
  \item \emph{Say it:} \textbf{â}, the vowel of \emph{father}
  \item \emph{Point to:} at the right edge of a Persian word, where your pen would begin
\end{itemize}

\begin{tcolorbox}[breakable,colback=teal!4,colframe=teal!35!black,title={Before you move on}]
How many strokes in alef? (\textbf{One}.)
\end{tcolorbox}
\section[lâm]{\fa{ل} — alef, with a bowl at its foot}

\label{lesson:FA-W15-lam}



\begin{quote}
Write \fa{ا}. One stroke, top to bottom.
\end{quote}

\subsection*{Script — the shape}
\begin{quote}
\fa{ل}
\end{quote}

Start exactly as you started alef — a tall stroke coming down — and then, instead of lifting, \textbf{swing left into a shallow bowl} at the foot and come back up a little.

One movement, no pen lift. If you can write \fa{ا} you are most of the way to \fa{ل} already; the difference is what happens at the bottom.

Its name is \textbf{lâm} and it says \textbf{l}, as in English \emph{long}.

\subsection*{Your turn}
\begin{itemize}
\raggedright
  \item \emph{Write it:} \fa{ا}, then \fa{ل} — the same descent, and then the bowl
  \item \emph{Say it:} \textbf{l}
\end{itemize}

\begin{tcolorbox}[breakable,colback=teal!4,colframe=teal!35!black,title={Before you move on}]
Do you lift the pen between the stroke and the bowl? (\textbf{No} — one movement.)
\end{tcolorbox}
\section[sin]{\fa{س} — three teeth, then a bowl}

\label{lesson:FA-W15-sin}



\begin{quote}
Write \fa{ل}. Descent, then bowl.
\end{quote}

\subsection*{Script — the shape}
\begin{quote}
\fa{س}
\end{quote}

Three small \textbf{teeth} — short upward points in a row — and then the line drops into a \textbf{bowl} and comes back up, the same kind of bowl you drew on lâm.

Remember you are working \textbf{right to left}: the first tooth you draw is the \textbf{rightmost} one.

All of it is one movement. No lifts at all.

Its name is \textbf{sin} and it says \textbf{s}, as in English \emph{see}.

\subsection*{Your turn}
\begin{itemize}
\raggedright
  \item \emph{Write it:} \fa{س} — three teeth from the right, then the bowl, without lifting
  \item \emph{Say it:} \textbf{s}
\end{itemize}

\begin{tcolorbox}[breakable,colback=teal!4,colframe=teal!35!black,title={Before you move on}]
Which tooth do you draw first? (The \textbf{rightmost}.)
\end{tcolorbox}
\section[mim]{\fa{م} — a loop, and a tail below the line}

\label{lesson:FA-W15-mim}



\begin{quote}
Write \fa{س}. Three teeth, then the bowl.
\end{quote}

\subsection*{Script — the shape}
\begin{quote}
\fa{م}
\end{quote}

A small \textbf{closed loop} sitting on the line, and from it a \textbf{tail} dropping straight down \textbf{below} the line.

The loop is small and round; keep it tight, or it starts to look like other letters you will meet later. The tail is what makes mim unmistakable standing on its own — very few letters drop below the line, and your eye will learn to use that. Joined into the middle of a word the tail is trimmed away and the loop is what survives, which is the same trimming you will meet everywhere.

One movement, no lift.

Its name is \textbf{mim} and it says \textbf{m}.

\subsection*{Your turn}
\begin{itemize}
\raggedright
  \item \emph{Write it:} \fa{م} — the small loop, then the tail dropping below the line
  \item \emph{Say it:} \textbf{m}
\end{itemize}

\begin{tcolorbox}[breakable,colback=teal!4,colframe=teal!35!black,title={Before you move on}]
Is the loop open or closed? (\textbf{Closed}.)
\end{tcolorbox}
\section[salâm]{Letters that hold hands, and two that will not}

\label{lesson:FA-W15-joining}



\begin{quote}
Write \fa{س}, then \fa{م}. Four letters are now yours.
\end{quote}

\subsection*{Script — the running line}
Persian is not written as separate letters sitting side by side. Letters \textbf{join}, running into one another along a single line, the way English cursive does — except that in Persian it is not a style. It is the ordinary way to write, and printed text does it too.

Because they join, most letters \textbf{change shape} depending on where they sit: a letter at the start of a joined run, one in the middle, one at the end, and one standing alone can all look different. They are the same letter.

That sounds like four times the work. It is not, and here is why: the changes are almost always \textbf{the same trimming}. A letter's identifying part stays, and the finishing flourish is cut short so the next letter can begin.

Look at sin, which you can already write:

\begin{quote}
alone: \textbf{\fa{س}} — in the middle of a run: \textbf{\fa{ـسـ}}
\end{quote}

The three teeth survive. The bowl is what gets trimmed away.

\subsection*{Script — and the letters that refuse}
Now the rule that saves you the most trouble.

\textbf{A few letters join to the letter before them, but never to the letter after.} They connect on the right, and refuse on the left, so the run \textbf{breaks} and the next letter starts fresh as though beginning a new word.

\textbf{alef is one of them.} So when you meet a gap in the middle of a Persian word, it is very often not a gap at all — it is one of these non-joining letters doing exactly what it always does.

That single fact turns a page of Persian from an unbroken puzzle into something with visible seams.

\begin{grammarlens}[title={What you've built}]
Four letters, and they spell a word you have been saying since the very first lesson:

\begin{quote}
\fa{سلام}
\end{quote}

Read it right to left: \textbf{\fa{س}} \emph{s}, \textbf{\fa{ل}} \emph{l}, \textbf{\fa{ا}} \emph{â}, \textbf{\fa{م}} \emph{m}. Say it: \textbf{salâm}.

Two things to notice. The short vowel between s and l is \textbf{not written} — Persian normally prints only the long vowels, and readers supply the rest. And the break before the final \fa{م} is alef refusing to join, exactly as promised.
\end{grammarlens}

\subsection*{Your turn}
\begin{itemize}
\raggedright
  \item \emph{Read it:} \fa{سلام} — name each letter from the right, then say the word
  \item \emph{Write it:} \fa{سلام} — and let the run break after alef
  \item \emph{Say it:} \textbf{salâm}
\end{itemize}

\begin{tcolorbox}[breakable,colback=teal!4,colframe=teal!35!black,title={Before you move on}]
Which part of sin survives when it joins? (The \textbf{teeth}; the bowl is trimmed.)
\end{tcolorbox}
\section[be]{\fa{ب} — a bowl, and one dot underneath}

\label{lesson:FA-W15-be}



\begin{quote}
Write \fa{سلام}. Four letters, one break.
\end{quote}

\subsection*{Script — the shape}
\begin{quote}
\fa{ب}
\end{quote}

A single shallow \textbf{bowl}, wider and flatter than the one on lâm, sitting on the line — and \textbf{one dot below it}.

Draw the bowl first, then lift and place the dot. That lift is the only one.

Its name is \textbf{be} and it says \textbf{b}.

Now the part worth more than the letter. \textbf{This bowl is a skeleton that several letters share.} On its own it means nothing; the \textbf{dots} decide which letter it is — how many, and whether they sit above or below.

That is why dots in Persian are not decoration and not optional. Move one dot from below to above and you have written a different letter. Your eye has to learn to look at the dots as hard as it looks at the shape.

\subsection*{Your turn}
\begin{itemize}
\raggedright
  \item \emph{Write it:} \fa{ب} — the shallow bowl, then one dot beneath it
  \item \emph{Say it:} \textbf{b}
\end{itemize}

\begin{tcolorbox}[breakable,colback=teal!4,colframe=teal!35!black,title={Before you move on}]
Is the bowl alone enough to identify a letter? (\textbf{No} — the dots decide.)
\end{tcolorbox}
\section[te, nun]{\fa{ت} and \fa{ن} — the same bowl, dotted differently}

\label{lesson:FA-W15-te-nun}



\begin{quote}
Write \fa{ب}. Bowl, then one dot below.
\end{quote}

\subsection*{Script — two dots above}
\begin{quote}
\fa{ت}
\end{quote}

\textbf{The same bowl as be.} Two dots, and they sit \textbf{above} it.

That is the entire difference, and it is a whole different letter: \textbf{te}, saying \textbf{t}.

\subsection*{Script — one dot above, and a deeper bowl}
\begin{quote}
\fa{ن}
\end{quote}

\textbf{One dot above} — and a bowl that is noticeably \textbf{deeper and rounder} than be's, curving well below the line before coming back up.

Its name is \textbf{nun} and it says \textbf{n}.

Standing alone, nun's deeper bowl helps you. In the middle of a joined run the bowl gets trimmed like everyone else's, and then the dot is doing nearly all the work: \textbf{one above} is nun, \textbf{two above} is te, \textbf{one below} is be.

\begin{grammarlens}[title={What you've built}]
Three letters, one shape:

\begin{quote}
\fa{ب} — one dot \textbf{below}
\end{quote}

>

\begin{quote}
\fa{ت} — two dots \textbf{above}
\end{quote}

>

\begin{quote}
\fa{ن} — one dot \textbf{above}
\end{quote}

You learned one bowl and got three letters out of it. That is the trade Persian keeps offering: fewer shapes than you feared, and dots you cannot afford to skim.
\end{grammarlens}

\subsection*{Your turn}
\begin{itemize}
\raggedright
  \item \emph{Write it:} \fa{ب}, then \fa{ت}, then \fa{ن} — same hand movement, different dots
  \item \emph{Say it:} \textbf{b}, \textbf{t}, \textbf{n}
\end{itemize}

\begin{tcolorbox}[breakable,colback=teal!4,colframe=teal!35!black,title={Before you move on}]
Which of the three has the deepest bowl? (\textbf{nun}.)
\end{tcolorbox}
\section[he]{\fa{ه} — the letter with four costumes}

\label{lesson:FA-W15-he}



\begin{quote}
Write \fa{ت}, then \fa{ن}. Same bowl, different dots.
\end{quote}

\subsection*{Script — the shape}
\begin{quote}
\fa{ه}
\end{quote}

A small \textbf{round} letter, standing on the line, no dots at all.

Its name is \textbf{he}. It says \textbf{h} at the start of a word — and at the \textbf{end} of a word it usually says nothing like an h: it carries a final \textbf{e} sound, the vowel that ends so many Persian words.

This letter will test what you learned about joining, because its four forms differ more than most:

\begin{quote}
alone: \textbf{\fa{ه}} — at the start: \textbf{\fa{هـ}} — in the middle: \textbf{\fa{ـهـ}} — at the end: \textbf{\fa{ـه}}
\end{quote}

Do not try to memorise four shapes. Recognise the \textbf{round body}, notice where in the run it is sitting, and let the shape follow from that. Persian readers are not holding four pictures in mind either.

\subsection*{Your turn}
\begin{itemize}
\raggedright
  \item \emph{Write it:} \fa{ه} — small and round, no dots
  \item \emph{Say it:} \textbf{h} at the front of a word; a soft \textbf{e} at the end of one
\end{itemize}

\begin{tcolorbox}[breakable,colback=teal!4,colframe=teal!35!black,title={Before you move on}]
Does he carry any dots? (\textbf{None}.)
\end{tcolorbox}
\section[vâv]{\fa{و} — one shape, three jobs, and no hand offered to the left}

\label{lesson:FA-W15-vav}



\begin{quote}
Write \fa{ه}. Small, round, no dots.
\end{quote}

\subsection*{Script — the shape}
\begin{quote}
\fa{و}
\end{quote}

A small \textbf{loop} at the top with a \textbf{tail} sweeping down and to the left, below the line. One movement.

Its name is \textbf{vâv}, and it is unusually hard-working. Depending on the word it is:

\begin{itemize}
\raggedright
  \item the consonant \textbf{v}
  \item the long vowel \textbf{u}, as in English \emph{food}
  \item the vowel \textbf{o}
\end{itemize}

You cannot tell which from the letter alone. You learn it with the word — and that is a fair trade, because it means Persian needs no separate letters for those vowels at all.

\textbf{And vâv is a non-joiner}, the second one this chapter has given you. Like alef, it takes the hand of the letter on its right and offers nothing to the left, so the run breaks after it.

\subsection*{Your turn}
\begin{itemize}
\raggedright
  \item \emph{Write it:} \fa{و} — the small loop, then the tail below the line
  \item \emph{Say it:} \textbf{v}, then \textbf{u}, then \textbf{o} — one letter, three jobs
\end{itemize}

\begin{tcolorbox}[breakable,colback=teal!4,colframe=teal!35!black,title={Before you move on}]
Can you tell from the shape alone whether vâv is v, u or o? (\textbf{No} — the word tells you.)
\end{tcolorbox}
\section[Practice]{Practice — seven words you already knew, now read}

\label{lesson:FA-C15-practice}



\begin{quote}
Write your nine letters in a row, right to left, saying each name as you go.
\end{quote}

\subsection*{Script — the words come back}
\textbf{Nothing new is taught here.} Every word below is one this book already gave you to say, and every letter in them is one you can now write.

Read each from the right, naming its letters before you say it:

\begin{quote}
\fa{سلام}
\end{quote}

>

\begin{quote}
\fa{بله}
\end{quote}

>

\begin{quote}
\fa{نه}
\end{quote}

>

\begin{quote}
\fa{ممنون}
\end{quote}

>

\begin{quote}
\fa{نان}
\end{quote}

>

\begin{quote}
\fa{نام}
\end{quote}

>

\begin{quote}
\fa{ماه}
\end{quote}

\emph{salâm} — \emph{bale} — \emph{na} — \emph{mamnun} — \emph{nân} — \emph{nâm} — \emph{mâh}.

Seven words. Nine letters. That ratio is the point: a small closed set of shapes unlocks far more than its own size, and it keeps doing so as the set grows.

\begin{culture}
Count the letters in \textbf{\fa{بله}} and count the sounds in \emph{bale}. Three letters, four sounds.

\textbf{Persian normally does not write its short vowels.} The long ones get letters — that is part of what alef and vâv are for — but the short \emph{a}, \emph{e} and \emph{o} are left off the page, and the reader supplies them.

This sounds alarming and is not, for the same reason English readers cope with \emph{read} and \emph{read}: you are not decoding letter by letter for long. You recognise the word.

It does mean one thing, though, and it is worth accepting early. \textbf{You cannot reliably sound out a Persian word you have never heard.} Reading and listening have to grow together in this language, more tightly than in one that spells every vowel. That is why this book taught you these seven words by ear first, and is only now handing you the page.
\end{culture}

\begin{grammarlens}[title={What you've built this chapter}]
Nine letters, and three ideas that matter more than any of them:

\begin{itemize}
\raggedright
  \item \textbf{the page runs right to left}, and your pen starts at the right edge
  \item \textbf{letters join}, changing shape by position — but the change is a trimming, not a new letter, and \textbf{alef and vâv refuse to join to the left}, which is what makes the gaps inside a word
  \item \textbf{dots carry the identity} — one bowl gave you \fa{ب}, \fa{ت} and \fa{ن}
\end{itemize}

Nine of Persian's thirty-two letters. The remaining twenty-three are mostly more of what you have just done: familiar skeletons, re-dotted.
\end{grammarlens}

\subsection*{Your turn}
\begin{itemize}
\raggedright
  \item \emph{Read it:} \fa{سلام} \fa{بله} \fa{نه} \fa{ممنون} — name the letters of each before saying it
  \item \emph{Write it:} \fa{نان} and \fa{نام} — one letter apart, and check which
  \item \emph{Say it:} \textbf{mamnun}, and count how many letters are written for it
\end{itemize}

\begin{tcolorbox}[breakable,colback=teal!4,colframe=teal!35!black,title={Before you move on}]
Which two of your nine letters refuse to join to the left? (\textbf{\fa{ا}} and \textbf{\fa{و}}.)
\end{tcolorbox}
